\chapter{Criterio de marco estable de evolución futura del corpus}

\section{Marco estable}
\begin{definition}
Se llama marco estable de evolución futura del corpus al dispositivo disciplinado mediante el cual MADELEINE puede organizar fases posteriores, expansiones laterales y desarrollos de largo alcance sin perder memoria editorial ni jerarquía interna.
\end{definition}

\begin{observation}
Su función es impedir que el porvenir del corpus se vuelva una masa indiferenciada de cierres parciales, reaperturas implícitas y crecimiento sin niveles.
\end{observation}

\section{Condición mínima de estabilidad evolutiva}
\begin{definition}
Existe marco estable de evolución futura cuando pueden declararse simultáneamente:
\begin{enumerate}
  \item una fase actual identificable;
  \item una diferencia legible entre fase cerrada, fase en transición y fase futura;
  \item una regla mínima de continuidad entre ciclos;
  \item una memoria editorial suficiente;
  \item una tipología mínima de expansión no contradictoria con la arquitectura previa.
\end{enumerate}
\end{definition}

\section{Tipología mínima}
\begin{definition}
El marco estable puede distinguir, al menos, entre:
\begin{enumerate}
  \item continuidad disciplinada;
  \item reapertura localizada;
  \item nueva fase editorial;
  \item expansión lateral;
  \item reintegración al corpus principal.
\end{enumerate}
\end{definition}

\begin{observation}
Esta tipología no agota todas las trayectorias futuras, pero basta para ordenar el crecimiento sin mezclar operaciones distintas bajo una sola palabra.
\end{observation}

\section{Estabilidad y apertura}
\begin{observation}
La estabilidad evolutiva no prohíbe expansión posterior. Solo exige que la expansión futura entre como fase reconocible, con memoria, trazabilidad y límites suficientemente explícitos.
\end{observation}

\section{Condición mínima de no estabilidad}
\begin{definition}
No debe declararse marco estable de evolución futura cuando la arquitectura propuesta dependa de ambigüedad sobre la fase actual, mezcla entre reapertura y nueva fase, o ausencia de memoria editorial suficiente.
\end{definition}

\section{Cláusula de alcance}
\begin{observation}
Este bloque compacta los criterios de continuidad entre ciclos, esquema general de fases y evolución futura de largo alcance. No constituye todavía una taxonomía exhaustiva del porvenir completo del sistema.
\end{observation}

